\section{Setting, scores and the evaluation problem}
\label{sec:method}

\paragraph{Compositional outputs.}
We call an output \emph{compositional} if it lies in the simplex $\simplex=\{p\in\mathbb{R}^K_{\ge0}:\sum_k p_k=1\}$.
Two settings are studied.
In Setting~A the \emph{target} is a composition, $Y\in\simplex$, and a method returns a region $C(X)\subset\simplex$.
In Setting~B the \emph{prediction} is a composition: a classifier returns $\hat p(x)\in\simplex$, the target is a label $Y\in\{1,\dots,K\}$, and a method returns a label set $C(x)\subset\{1,\dots,K\}$.

\paragraph{Split conformal prediction.}
Let a nonconformity score $s(x,y)$ be fitted on a training split, let $S_1,\dots,S_n$ be the scores on a calibration split of size $n$, and let $S_{n+1}$ be the score of a test pair.
For miscoverage level $\alpha$ set $k=\lceil (n+1)(1-\alpha)\rceil$ and $\ell=n+1-k=\lfloor (n+1)\alpha\rfloor$.
The threshold is $\hat q=S_{(k)}$ (the $k$-th smallest calibration score; $\hat q=\infty$ if $k>n$) and $C(x)=\{y:s(x,y)\le \hat q\}$.

\subsection{The exact law of coverage}
\label{sec:theory}
Define the conditional coverage given the calibration data $V=\Pbb\{S_{n+1}\le\hat q\mid S_1,\dots,S_n\}$.

\begin{proposition}[coverage law]\label{prop:beta}
If $S_1,\dots,S_{n+1}$ are exchangeable with a continuous joint law (no ties almost surely), then $V\sim\mathrm{Beta}(n+1-\ell,\ \ell)$.
Consequently (i) $\Ebb V=(n+1-\ell)/(n+1)\in[1-\alpha,\,1-\alpha+\tfrac1{n+1})$; (ii) $\operatorname{Var}V=\dfrac{(n+1-\ell)\,\ell}{(n+1)^2(n+2)}\approx\dfrac{\alpha(1-\alpha)}{n+2}$; and (iii) the number of covered points among $m$ fresh test points, given the calibration data, is $\mathrm{Bin}(m,V)$, so unconditionally it is $\mathrm{BetaBinomial}(m,\,n+1-\ell,\,\ell)$.
This law is the same for every score function $s$, hence for every choice of geometry on $\simplex$.
\end{proposition}

The law is a known property of split conformal prediction \citep{vovk2012conditional,angelopoulos2023gentle}; we state it because the rest of the paper uses it as the reference against which every score, geometry and estimator is checked.

\begin{proof}
For i.i.d.\ scores with continuous cdf $F$, $U_i=F(S_i)$ are i.i.d.\ uniform and $V=F(S_{(k)})=U_{(k)}$ is the $k$-th order statistic of $n$ uniforms, which is $\mathrm{Beta}(k,n+1-k)=\mathrm{Beta}(n+1-\ell,\ell)$.
For merely exchangeable scores with no ties, the same holds conditionally on the order statistics by the usual randomisation argument (the rank of $S_{n+1}$ among the $n+1$ scores is uniform), and $V$ depends only on ranks.
The mean and variance are those of the Beta law, and (iii) follows because test points are, given the calibration data, i.i.d.\ Bernoulli$(V)$ for i.i.d.\ data.
\end{proof}

\begin{remark}[ties]
If the score has atoms (rounded compositions, zeros replaced by a constant, discrete $\hat p$), coverage remains $\ge1-\alpha$ but the upper bound in (i) is lost, and coverage is conservative by a score-dependent amount (tested in Section~\ref{sec:experiments}).
\end{remark}

\paragraph{What this implies for evaluation.}
(a) Marginal coverage cannot separate score geometries, which differ only through region size and conditional coverage.
(b) The spread of the \emph{measured} coverage is dominated by the calibration draw, $\operatorname{sd}\approx\sqrt{\alpha(1-\alpha)/(n+2)}$, which a larger test set does not remove.
(c) Certifying $|V-(1-\alpha)|\le\varepsilon$ with probability $1-\delta$ requires the $n$ with $\Pbb(|V-(1-\alpha)|\le\varepsilon)\ge1-\delta$, computed from the Beta cdf; the normal approximation gives $n\gtrsim\alpha(1-\alpha)z_{1-\delta/2}^2/\varepsilon^2$.

\subsection{Bias of worst-cell coverage estimators}
\label{sec:minbias}
Conditional coverage is evaluated by partitioning the feature space (or the set-size range, or the label set) into $G$ cells and reporting the worst cell.
Condition on the fitted model and calibration data, let $c_g$ be the true coverage of cell $g$, and let $m_g$ test points fall in it.
The cell estimates are $\hat c_g=B_g/m_g$ with $B_g\sim\mathrm{Bin}(m_g,c_g)$ independent across cells (disjoint test points).

\begin{proposition}[worst-cell estimator]\label{prop:min}
For equal cell sizes $m_g=m$,
\[
\Ebb\Big[\min_g \hat c_g\Big]=\frac1m\sum_{j=1}^{m}\prod_{g=1}^{G}\Pbb\{\mathrm{Bin}(m,c_g)\ge j\},
\]
which is at most $\min_g c_g$, with strict inequality whenever some $\hat c_g$ is non-degenerate.
If all $c_g=c$ and $m$ is large, $\Ebb[\min_g\hat c_g]\approx c-\sqrt{c(1-c)/m}\;\mu_G$ with $\mu_G=\Ebb\max_{g\le G}Z_g$ for i.i.d.\ standard normals ($\mu_{10}\approx1.54$).
\end{proposition}
\begin{proof}
$\min_g\hat c_g$ takes values in $\{0,1/m,\dots,1\}$, so $\Ebb\min=\sum_{j=1}^m\frac1m\Pbb\{\min_g B_g\ge j\}$ and independence gives the product.
The inequality is $\Ebb\min_g\hat c_g\le\min_g\Ebb\hat c_g=\min_g c_g$ (Jensen for the concave $\min$).
The Gaussian approximation is the expected maximum of $G$ i.i.d.\ normals applied to $-(\hat c_g-c)/\sqrt{c(1-c)/m}$.
\end{proof}

The bias grows with $G$ and shrinks with $m=n_{\rm te}/G$ and is of the order of the differences between methods one hopes to resolve; rankings are affected through differences in the spread of $c_g$ across methods, not through a common offset.
